\scanpage{595}

\mspage{4}

Soit alors $A$ l'ens des $\unsure{\xi|\zeta}\in E$ tels que $g(\xi)\ne0$, $B=E\setminus A$ l'ens des $\xi\in E$ tels que $g(\xi)\unsure{\ne|=}0$. \inserted{On a donc $B\notin\Phi$ \deleted{\unsure{parce que|puisque}} puisque $f\ne0$, et} comme $\Phi$ est un ultrafiltre, on a donc $A\in\Phi$.

Mais soit $g'\inserted{\in k_E}$ défini par $g'(\xi)=f(\xi)^{-1}$ si $\xi\in A$, \inserted{$g'(\xi)=1$ si $\xi\in B$,} et soit $f'$ l'image de $g'$ dans $k_\Phi$. On voit alors qu'on a

\[
ff'=f'f=1\ ,
\]

ce qui montre bien que $k_\Phi$ est un corps.

Supposons maintenant que tous les $k_\xi$ $(\xi\in E)$ soient égaux à un $=$ corps $k$. Alors les $k_J$ \unsure{forment} un système projectif de $k$-algèbres, les $k_\Phi$ un syst. inductif de $k$-algèbres. Si $k$ est commutatif, il en est de $=$ des $k_J$, $k_\Phi$, et en particulier \unsure{pour} $\Phi$ un ultrafiltre, $k_\Phi$ est une \uline{extension} du corps $k$.

Revenons au cas général, \inserted{mais les $k_\xi$ comm.} \deleted{je dis que} \inserted{Soit} \deleted{\uline{Proposition}} $\unsure{\mathscr F|\mathcal F}_0\subset\mathscr F$ l'ens des ultrafiltres sur $E$,

Notons que pour $\unsure{\forall|\mathrm{tt}}\Phi\in\mathscr F$, $k_E\mathrel{\unsure{-|\to}}k_\Phi$ est surjectif. Donc si $\Phi\in\mathscr F_{\unsure{0}}$, $k_\Phi$ s'identifie à un corps quotient de l'\unsure{anneau} $k_E$.

% archive sheet number: 594
